% --- 5_1_trong-so-N10.png ---
\begin{frame}{Từ điểm importance sang trọng số xoá: ví dụ $N=10$}
\begin{columns}[T]
\begin{column}{0.52\linewidth}
\footnotesize
\begin{itemize}
    \item Kịch bản giả định $N=10$ Gaussian với điểm Pruning\_i cho trước:
    \[
    P = [0{,}0.05,0{,}2,0{,}4,0{,}6,0{,}8,0{,}9,0{,}95,0{,}99,1{,}0]
    \]
    \item Tập ứng viên $C=\{3,6,8,9,10\}$ (chỉ số 1-based), ngân sách xoá
    \[
    \text{budget} = \lfloor 0{,}5\cdot |C| \rfloor = \lfloor 0{,}5 \cdot 5 \rfloor = 2
    \]
    \item Trọng số bị xoá tỉ lệ nghịch với $(1-P_i)$:
    \[
    w_i = \frac{1}{10^{-6} + (1-P_i)}
    \]
    \item Chuẩn hoá: $p_i = w_i / \sum_j w_j$.
    \item $w_i$ \textbf{phân kỳ} khi $P_i \to 1$: Gaussian số 10 ($P_{10}=1{,}0$) chiếm tới $99{,}99\%$ xác suất bị rút ở lượt đầu tiên -- gần như chắc chắn bị xoá.
\end{itemize}
\end{column}
\begin{column}{0.48\linewidth}
\centering
\includegraphics[width=0.98\linewidth]{5_1_trong-so-N10.png}
\captionof{figure}{\footnotesize $P_i$, trọng số $w_i$ (log) và xác suất $p_i$ (log) của 10 Gaussian; đỏ = thuộc tập ứng viên $C$.}
\end{column}
\end{columns}
\end{frame}

% --- 5_2_mo-phong-2000-lan.png ---
\begin{frame}{Kiểm chứng bằng mô phỏng: 2000 lần lấy mẫu multinomial}
\footnotesize
\begin{itemize}
    \item Lấy mẫu \texttt{multinomial\_no\_replacement}$(w,\text{budget}{=}2)$ lặp lại $T=2000$ lần (không hoàn lại, giống \texttt{torch.multinomial}).
    \item Xác suất lý thuyết $P(i \in S)$ với budget $=2$ tính chính xác bằng công thức bao gồm cả cặp:
    \[
    P(i \in S) = p_i + \sum_{j \neq i} p_j \cdot \frac{w_i}{W - w_j}, \qquad p_i = \frac{w_i}{W}, \; W=\sum_k w_k
    \]
    \item Tần suất thực nghiệm $i \in S$ qua 2000 lần khớp sát với $P(i\in S)$ lý thuyết -- xác nhận cơ chế lấy mẫu đúng như thiết kế.
    \item Tần suất \emph{thực sự bị xoá} = $S \cap C$: một Gaussian có thể được rút vào $S$ nhưng không thuộc $C$ nên không bị xoá (ví dụ G7 hay được rút nhưng không nằm trong tập ứng viên).
    \item Số Gaussian bị xoá thực tế trung bình $\approx$ giá trị hiển thị trên biểu đồ (thấp hơn budget=2 vì $S\cap C$ có thể nhỏ hơn $S$).
\end{itemize}
\centering
\includegraphics[width=0.62\linewidth]{5_2_mo-phong-2000-lan.png}
\captionof{figure}{\footnotesize Trái: $P(i\in S)$ lý thuyết vs tần suất thực nghiệm. Giữa: xác suất/tần suất bị xoá $S\cap C$. Phải: phân bố số Gaussian bị xoá thực tế mỗi lượt.}
\end{frame}

% --- 5_3_topk-vs-multinomial.png ---
\begin{frame}{Top-k cứng vs.\ Multinomial: hai chiến lược pruning}
\begin{columns}[T]
\begin{column}{0.5\linewidth}
\footnotesize
\begin{itemize}
    \item Mô phỏng $N=300$ Gaussian, điểm số $P_i$ chuẩn hoá min--max từ \texttt{Beta}$(2{,}2{,}5)$; tập ứng viên $C$ chọn ngẫu nhiên $40\%$ quần thể (độc lập với $P_i$), budget $=\lfloor 0{,}5|C|\rfloor$.
    \item \textbf{Top-k cứng (deterministic):} sắp xếp $C$ giảm dần theo $P_i$, xoá đúng \texttt{budget} phần tử điểm thấp nhất.
    \begin{itemize}
        \item[+] Dễ đoán, luôn xoá đúng số lượng.
        \item[--] Cắt sạch theo một ngưỡng cứng -- có thể xoá nhầm cả một cụm liền kề có điểm số tương tự nhau.
    \end{itemize}
    \item \textbf{Multinomial (stochastic):} rút mẫu không hoàn lại theo $w_i \propto 1/(10^{-6}+1-P_i)$ trên toàn $N$, giao với $C$.
    \begin{itemize}
        \item[+] Đa dạng hơn, tránh xoá đồng loạt một vùng.
        \item[--] Có nhiễu: số lượng xoá thực tế thường \emph{ít hơn} budget vì $S\cap C$ nhỏ hơn $S$.
    \end{itemize}
\end{itemize}
\end{column}
\begin{column}{0.5\linewidth}
\centering
\includegraphics[width=0.95\linewidth]{5_3_topk-vs-multinomial.png}
\captionof{figure}{\footnotesize (a) top-k cứng cắt theo một đường ngưỡng; (b) multinomial xoá ít hơn budget, có ngẫu nhiên; (c) phân bố điểm còn lại; (d) xác suất bị xoá tăng mạnh khi $P_i\to 1$.}
\end{column}
\end{columns}
\end{frame}

% --- 5_4_tranh-xoa-ca-cum.png ---
\begin{frame}{Vì sao lấy mẫu ngẫu nhiên tránh xoá sạch một cụm}
\footnotesize
\begin{itemize}
    \item Kịch bản: $N=300$ điểm 2D, một \textbf{cụm ``nóng''} gần tâm $(7{,}5,7{,}5)$ bán kính $1{,}6$ có $P_i$ đồng đều cao $\in[0{,}86;0{,}95]$; phần còn lại $P_i\in[0{,}55;0{,}92]$.
    \item Budget xoá $=15\%\cdot N = 45$ Gaussian, tất cả đều là ứng viên ($C = $ toàn bộ quần thể).
    \item \textbf{Top-k cứng:} vì mọi điểm trong cụm đều có $P_i$ cao và tương tự nhau, thuật toán xoá gần như \emph{toàn bộ} cụm cùng lúc $\Rightarrow$ tạo một \textbf{lỗ hổng hình học} tại vị trí đó.
    \item \textbf{Multinomial:} do tính ngẫu nhiên trong lấy mẫu theo trọng số, chỉ một phần cụm bị xoá ở mỗi lượt, phần còn lại của cụm vẫn được giữ lại để tái đánh giá ở lượt sau $\Rightarrow$ tránh mất mát hình học đột ngột.
    \item Đây là lý do chính khiến pruning dùng lấy mẫu xác suất thay vì cắt ngưỡng cứng tuyệt đối.
\end{itemize}
\centering
\includegraphics[width=0.68\linewidth]{5_4_tranh-xoa-ca-cum.png}
\captionof{figure}{\footnotesize Trái: bản đồ $P_i$ theo vị trí (vòng đỏ = cụm điểm cao). Giữa/phải: các điểm bị xoá bởi top-k cứng (xoá gần hết cụm) vs multinomial (rải rác hơn).}
\end{frame}

% --- 5_5_sigmoid-inverse.png ---
\begin{frame}{Sigmoid nghịch đảo (logit) để reset opacity}
\footnotesize
\begin{itemize}
    \item Tham số thực sự được tối ưu không phải $\alpha\in(0,1)$ mà là \texttt{\_opacity} (logit), với
    \[
    \alpha = \sigma(\text{logit}) = \frac{1}{1+e^{-\text{logit}}}, \qquad
    \text{logit} = \sigma^{-1}(\alpha) = \log\frac{\alpha}{1-\alpha}
    \]
    \item Khi cần ép/reset opacity về một giá trị nhỏ $\alpha_{\text{target}}$, ta phải gán lại \emph{logit}:
    \[
    \text{logit}_{\text{new}} = \sigma^{-1}(\min(\alpha,\alpha_{\text{target}}))
    \]
    vì gán trực tiếp lên $\alpha$ sẽ phá vỡ ràng buộc $(0,1)$ và làm sai gradient.
    \item Các mốc quan trọng: prune $\alpha=0{,}005$; reset $\alpha=0{,}01$; final prune $\alpha=0{,}1$; clamp $\alpha=0{,}8$.
    \item Bước $0{,}9\to 0{,}8$: logit $2{,}197 \to 1{,}386$ (bước nhỏ). Bước $0{,}9\to 0{,}01$: logit $2{,}197 \to -4{,}595$ (bước rất xa) $\Rightarrow$ reset là một cú "giật lùi" mạnh trong không gian logit.
\end{itemize}
\centering
\includegraphics[width=0.68\linewidth]{5_5_sigmoid-inverse.png}
\captionof{figure}{\footnotesize Trái: hàm sigmoid với các mốc opacity đánh dấu trên trục logit. Phải: hàm nghịch đảo $\sigma^{-1}(\alpha)$ trên thang log của $\alpha$.}
\end{frame}

% --- 5_6_opacity-theo-iteration.png ---
\begin{frame}{Quỹ đạo opacity qua thời gian: tăng rồi bị reset định kỳ}
\footnotesize
\begin{itemize}
    \item Mô phỏng minh hoạ 4 Gaussian với "sức kéo" gradient khác nhau lên logit mỗi vòng: hữu ích ($k{=}+0{,}006$), trung bình ($k{=}+0{,}002$), vô dụng ($k{=}-0{,}001$), hồi sinh muộn ($k{=}0$ rồi $+0{,}004$ sau $t{=}7000$).
    \item Mỗi 100 vòng (trong khoảng $500<t<15000$), logit bị clamp: $\text{logit}\leftarrow\min(\text{logit}, \sigma^{-1}(0{,}8))$ -- ngăn opacity vượt $0{,}8$ ngay sau densify.
    \item Mỗi \texttt{opacity\_reset\_interval} $=3000$ vòng: $\text{logit}\leftarrow\min(\text{logit}, \sigma^{-1}(0{,}01))$ -- kéo \emph{toàn bộ} opacity về gần $0$, bất kể trước đó cao bao nhiêu.
    \item Sau reset, Gaussian hữu ích tự "leo" trở lại nhanh (gradient dương mạnh); Gaussian vô dụng rơi xuống dưới ngưỡng prune $\alpha=0{,}005$ và bị xoá ở lần prune multinomial kế tiếp.
    \item Cơ chế này cho các Gaussian "được nghi ngờ" một cơ hội chứng minh lại giá trị, thay vì xoá vĩnh viễn ngay lập tức.
\end{itemize}
\centering
\includegraphics[width=0.72\linewidth]{5_6_opacity-theo-iteration.png}
\captionof{figure}{\footnotesize Quỹ đạo $\alpha_i(t)$ (log) của 4 Gaussian điển hình; các đường đứt nét đánh dấu mốc reset mỗi 3000 vòng, ngưỡng clamp 0,8 và ngưỡng prune 0,005.}
\end{frame}

% --- 5_7_histogram-truoc-sau-reset.png ---
\begin{frame}{Histogram phân bố opacity: trước và sau một lần reset}
\footnotesize
\begin{itemize}
    \item Mô phỏng $N=20000$ Gaussian: $60\%$ "hữu ích" có $\alpha_0\sim\text{Beta}(6,2)$ (thiên về giá trị cao), $40\%$ "lơ lửng" có $\alpha_0\sim\text{Beta}(1{,}5,4)$ (thiên về giá trị thấp).
    \item \textbf{Trước reset:} phân bố $\alpha_0$ đa dạng, trải rộng trên toàn khoảng $(0,1)$.
    \item \textbf{Ngay sau reset:} áp dụng $\text{logit}_1 = \min(\text{logit}_0, \sigma^{-1}(0{,}01))$ cho \emph{mọi} Gaussian $\Rightarrow$ toàn bộ phân bố dồn về sát $\alpha\approx 0{,}01$, gần như đồng nhất, mất hết sự đa dạng trước đó.
    \item \textbf{$\sim$500 vòng sau (minh hoạ):} nhóm hữu ích leo lại nhanh (gain trung bình $+3{,}5$ logit), nhóm vô dụng leo chậm hoặc đứng yên (gain trung bình $-0{,}3$ logit) $\Rightarrow$ phân bố tách lại thành hai nhóm.
    \item Tỉ lệ $\alpha<0{,}005$ (dưới ngưỡng prune) được ghi trực tiếp trên từng histogram để so sánh ba giai đoạn.
\end{itemize}
\centering
\includegraphics[width=0.85\linewidth]{5_7_histogram-truoc-sau-reset.png}
\captionof{figure}{\footnotesize Histogram $\alpha$ (thang log) trước reset, ngay sau reset, và $\sim$500 vòng sau reset; đường đứt = ngưỡng prune 0,005 và 0,01.}
\end{frame}

% --- 5_8_adam-sau-reset.png ---
\begin{frame}{Trạng thái Adam sau reset: có nên xoá luôn moment?}
\footnotesize
\begin{itemize}
    \item Adam cập nhật theo $m_t=\beta_1 m_{t-1}+(1-\beta_1)g$, $v_t=\beta_2 v_{t-1}+(1-\beta_2)g^2$, với bias-correction $\hat m_t = m_t/(1-\beta_1^t)$, $\hat v_t=v_t/(1-\beta_2^t)$, bước cập nhật $=\text{lr}\cdot \hat m_t/(\sqrt{\hat v_t}+\epsilon)$.
    \item Code thật (\texttt{replace\_tensor\_to\_optimizer}) khi reset opacity: đặt lại $(m,v)=(0,0)$ cho nhóm opacity, nhưng \textbf{không reset bộ đếm bước $t$ (step) của Adam} -- step vẫn tiếp tục từ giá trị cũ (vd. $t{=}6000$).
    \item Hệ quả: với $t$ lớn, $\beta_1^t\approx\beta_2^t\approx 0$ nên bias-correction gần như \emph{không còn hiệu lực đúng nghĩa} (được thiết kế cho $t$ nhỏ) $\Rightarrow$ bước cập nhật đầu tiên sau reset lớn bất thường, xấp xỉ
    \[
    \frac{\hat m_1}{\sqrt{\hat v_1}} \approx \frac{(1-\beta_1)g}{\sqrt{(1-\beta_2)g^2}} = \frac{0{,}1}{\sqrt{0{,}001}} \approx 3{,}16 \times \text{lr}
    \]
    \item So với Adam khởi tạo hoàn toàn mới (step$=0$, bias-correction đúng), phiên bản "giữ step" tạo một cú nhảy tham số lớn hơn nhiều ngay sau reset, rồi hội tụ dần về $1\times$lr ở trạng thái ổn định.
    \item Nếu không reset cả state Adam: optimizer "nhớ" quán tính cũ từ trước reset, có thể khiến opacity hồi phục lệch hướng hoặc chậm hơn dự kiến so với một khởi động sạch.
\end{itemize}
\centering
\includegraphics[width=0.6\linewidth]{5_8_adam-sau-reset.png}
\captionof{figure}{\footnotesize So sánh độ lớn bước cập nhật (chuẩn hoá theo lr) của Adam "giữ step cũ, reset $(m,v){=}(0,0)$" (code thật) vs Adam khởi tạo hoàn toàn mới.}
\end{frame}
